% Part III-C, results of the offline prefix study (included by main_en.tex).

\section{Can a judge tell a promising prefix?}\label{sec:res-prefix}

The base policy's \NUM{ps_problems}$\times$4 solutions gave \NUM{ps_prefixes} prefixes and \NUM{ps_continuations}
verified continuations, of which \NUM{ps_contacc} are correct. A prefix's value spreads over the whole range. At
128 tokens the mean value is \NUM{ps_value_128} and \NUM{ps_dead_128} of the prefixes are dead (no correct
continuation). At 512 tokens these are \NUM{ps_value_512} and \NUM{ps_dead_512}, because the solutions that are still
running by then are the harder ones. So there is something to choose: among one problem's prefixes, some are worth
continuing and some are not.

\paragraph{All signals carry information, none separates a problem's prefixes well.} Table~\ref{tab:prefix-auc} gives
each signal's AUROC for a continuation ending correct. Pooled over problems, every signal is clearly informative
(\NUM{auc_128_logprob}--\NUM{auc_128_artrack} at 128 tokens). Most of that is the difficulty of the problem, which every
signal partly sees. Within problems, the comparison that matters for choosing which prefixes to continue, the AUROCs
fall to \NUM{aucw_128_kevcorrect}--\NUM{aucw_128_logprob} at 128 tokens and reach at most \NUM{aucw_512_artrack} at
512 tokens. \kev{}'s \emph{on track} question does not do better than the AR judge of the same size asked the same
question (within problems \NUM{aucw_256_kevtrack} against \NUM{aucw_256_artrack} at 256 tokens, a difference of
\NUM{psdw_256_kevtrackartrack} \ci{\NUM{psdwlo_256_kevtrackartrack}}{\NUM{psdwhi_256_kevtrackartrack}}; 95\% intervals
from resampling problems). Nor does it beat the policy's own mean token log-probability, which costs nothing
(\NUM{aucw_256_logprob}; \kev{} minus log-probability \NUM{psdw_256_kevtracklogprob}
\ci{\NUM{psdwlo_256_kevtracklogprob}}{\NUM{psdwhi_256_kevtracklogprob}}). Pooled over problems both judges are ahead of
the log-probability (\NUM{psd_256_kevtracklogprob} \ci{\NUM{psdlo_256_kevtracklogprob}}{\NUM{psdhi_256_kevtracklogprob}}
for \kev{}), so what the judges add is mostly a sense of how hard the problem is, which does not help to choose among
one problem's prefixes. Asking \kev{} Part II's \emph{correct} question about an unfinished solution works worse than
asking whether it is on track.

\begin{table}[t]
\centering\small
\begin{tabular}{@{}lccccccccc@{}}
\toprule
 & \multicolumn{3}{c}{128 tokens} & \multicolumn{3}{c}{256 tokens} & \multicolumn{3}{c}{512 tokens} \\ \cmidrule(lr){2-4}\cmidrule(lr){5-7}\cmidrule(lr){8-10}
signal & AUROC & within & $\rho$ & AUROC & within & $\rho$ & AUROC & within & $\rho$ \\ \midrule
\kev{} \emph{on track} & 0.76 & 0.61 & 0.53 & 0.77 & 0.64 & 0.52 & 0.77 & 0.65 & 0.46 \\
\kev{} \emph{correct} & 0.70 & 0.59 & 0.40 & 0.75 & 0.59 & 0.48 & 0.77 & 0.58 & 0.47 \\
AR \emph{on track} & 0.77 & 0.63 & 0.55 & 0.79 & 0.64 & 0.56 & 0.79 & 0.67 & 0.50 \\
policy log-prob. & 0.67 & 0.64 & 0.34 & 0.72 & 0.66 & 0.42 & 0.73 & 0.64 & 0.39 \\
\bottomrule
\end{tabular}

\caption{How well each signal predicts that a continuation of the prefix ends correct, by prefix length. AUROC pooled
over all prefixes and within problems (pairs of continuations of the same problem with different outcomes), and the
Spearman correlation $\rho$ with the prefix's Monte-Carlo value. 500 validation problems, 4 solutions each, 4
continuations per prefix.}
\label{tab:prefix-auc}
\end{table}

\paragraph{Selecting prefixes helps a little, and any signal will do.} On the problems whose four solutions all run past
the cut, keeping the two prefixes a signal ranks highest raises the mean value of what is continued
(Table~\ref{tab:prefix-selection}). At 128 tokens \kev{}'s \emph{on track} gains \NUM{gain_128_kevtrack}
\ci{\NUM{gainlo_128_kevtrack}}{\NUM{gainhi_128_kevtrack}} over the \NUM{ps_valueall_128} of keeping all four. The AR
judge gains \NUM{gain_128_artrack} and the policy's log-probability \NUM{gain_128_logprob}. The best gains are at
256 tokens, \NUM{gain_256_kevtrack} for \kev{}, \NUM{gain_256_artrack} for the AR judge and \NUM{gain_256_logprob} for
the log-probability. The oracle, which ranks by two of the four continuations and is scored on the other two, gains
\NUM{gain_256_oracle} there. With four continuations per prefix the value itself is a noisy estimate, so the oracle's
gain underestimates what perfect knowledge would give. At 128 tokens the signals gain about as much as this oracle, at
256 tokens about two thirds of it. In absolute terms the gains are small: a few more correct continuations in a
hundred.

\begin{table}[t]
\centering\small
\begin{tabular}{@{}lccc@{}}
\toprule
keep 2 of 4 by & 128 tokens & 256 tokens & 512 tokens \\ \midrule
problems & 401 & 253 & 93 \\
value of all 4 & 0.386 & 0.324 & 0.165 \\ \midrule
\kev{} \emph{on track} & $+$0.046 \ci{0.034}{0.058} & $+$0.064 \ci{0.047}{0.082} & $+$0.040 \ci{0.015}{0.067} \\
\kev{} \emph{correct} & $+$0.040 \ci{0.028}{0.052} & $+$0.043 \ci{0.026}{0.060} & $+$0.021 \ci{$-$0.003}{0.046} \\
AR \emph{on track} & $+$0.045 \ci{0.033}{0.057} & $+$0.069 \ci{0.053}{0.085} & $+$0.048 \ci{0.022}{0.075} \\
policy log-prob. & $+$0.052 \ci{0.041}{0.063} & $+$0.065 \ci{0.048}{0.083} & $+$0.040 \ci{0.014}{0.067} \\
oracle (split half) & $+$0.048 \ci{0.033}{0.063} & $+$0.098 \ci{0.080}{0.117} & $+$0.059 \ci{0.034}{0.086} \\
\bottomrule
\end{tabular}

\caption{Mean value of the 2 prefixes (of 4) a signal keeps, minus the mean value of all 4, on the problems whose 4
solutions all run past the cut. 95\% bootstrap intervals over problems. The oracle ranks by half of the continuations
and is scored on the other half.}
\label{tab:prefix-selection}
\end{table}

\paragraph{A fixed pruning threshold rarely fires.} The rule prunes a prefix when the judge's
$p(\text{yes})$ is below 0.3 (Table~\ref{tab:prefix-threshold}). \kev{}'s \emph{on track} probability is almost never that
low. The rule prunes \NUM{thrpruned_128_kevtrack} of the prefixes at 128 tokens. \NUM{thrdead_128_kevtrack} of those are
dead, but they are only \NUM{thrcaught_128_kevtrack} of all dead prefixes. With the \emph{correct} question the rule
prunes more (\NUM{thrpruned_128_kevcorrect}), with fewer dead prefixes among the pruned (\NUM{thrdead_128_kevcorrect}),
and it throws away \NUM{thrlost_128_kevcorrect} of all correct continuations. A fixed threshold on an uncalibrated
probability is not a usable pruning rule here; ranking within a problem is.

\begin{table}[t]
\centering\small
\begin{tabular}{@{}lccccccccc@{}}
\toprule
 & \multicolumn{3}{c}{128 tokens} & \multicolumn{3}{c}{256 tokens} & \multicolumn{3}{c}{512 tokens} \\ \cmidrule(lr){2-4}\cmidrule(lr){5-7}\cmidrule(lr){8-10}
signal & pruned & dead & lost & pruned & dead & lost & pruned & dead & lost \\ \midrule
\kev{} \emph{on track} & 1.6\% & 84\% & 0.2\% & 2.1\% & 85\% & 0.3\% & 1.5\% & 87\% & 0.3\% \\
\kev{} \emph{correct} & 12.1\% & 61\% & 6.4\% & 12.1\% & 72\% & 4.2\% & 17.1\% & 76\% & 8.2\% \\
AR \emph{on track} & 4.6\% & 77\% & 1.2\% & 8.0\% & 78\% & 1.7\% & 14.4\% & 87\% & 3.9\% \\
\bottomrule
\end{tabular}

\caption{The threshold rule, prune when $p(\text{yes})<0.3$: the fraction of prefixes pruned, the fraction of the pruned that
were dead, and the fraction of all correct continuations that the pruned prefixes would have produced.}
\label{tab:prefix-threshold}
\end{table}
